\section{주분석의 실증전략}

\subsection{공동 주결과의 두 패널}

Panel A는 삼쌍승과 비교 승식의 배당이 모두 상한에 걸리지 않은 3,321개 경주에서
점가격을 비교한다. Panel B는 완전 지지집합을 가진 전체 19,284개 경주에서
게시 배당의 검열·반올림 규칙과 양립하는 가격집합을 사용한다. 경주 $r$과 승식
$m$에서 이 집합을 $\mathcal P^m_r$라 하자. 상한에 걸린 총지급배율은 실제
배율의 하한만 알려 주므로 원시 역배당률의 상한을 주고, 반올림된 비상한 배당은
원시 역배당률의 양쪽 경계를 준다. 이 성분별 구간을 같은 풀 안에서 정규화하여
$\mathcal P^m_r$를 구성한다.

손실함수 $L$의 전체표본 식별구간은
\[
  \left[
    \inf_{p\in\mathcal P^m_r,\,q\in\mathcal P^T_r}
      L(p,A^m_rq),\quad
    \sup_{p\in\mathcal P^m_r,\,q\in\mathcal P^T_r}
      L(p,A^m_rq)
  \right]
\]
로 계산한다. 기준모형과의 손실 차이는 각 손실의 별도 경계를 뺄셈하지 않고,
동일한 허용 가격벡터 위에서 손실 차이 자체를 최소화·최대화한다. 따라서 공통
가격 불확실성이 중복 계산되지 않는다. 표본불확실성은 경주 단위 부트스트랩으로
식별구간의 양 끝점에 반영한다. $\epsilon$이 필요한 로그손실과 calibration에는
Panel A와 Panel B에서 같은 사전 고정값을 사용한다.

\subsection{주요 평가량}

승식 $m$과 경주 $r$에서 실제 정규화 가격 $p^m_r$와 삼쌍승 재구성 가격
$\widehat p^{m\leftarrow\T}_r$의 차이를 다음과 같이 측정한다.
\begin{align*}
  \operatorname{TV}_{r,m}
    &=\frac12\sum_c\left|p^m_{r,c}-
      \widehat p^{m\leftarrow\T}_{r,c}\right|,\\
  \operatorname{RMSLE}_{r,m}
    &=\left[\frac{1}{C_{r,m}}\sum_c
      \left\{\log(p^m_{r,c}+\epsilon)-
      \log(\widehat p^{m\leftarrow\T}_{r,c}+\epsilon)\right\}^2
      \right]^{1/2},\\
  \operatorname{JS}_{r,m}
    &=\frac12\sum_c p^m_{r,c}\log\frac{p^m_{r,c}}{s^m_{r,c}}
      +\frac12\sum_c \widehat p^{m\leftarrow\T}_{r,c}
       \log\frac{\widehat p^{m\leftarrow\T}_{r,c}}{s^m_{r,c}},\\
  s^m_{r,c}
    &=\frac12\left(p^m_{r,c}+\widehat p^{m\leftarrow\T}_{r,c}\right),\\
  \operatorname{MAE}_{r,m}
    &=\frac{1}{C_{r,m}}\sum_c
      \left|p^m_{r,c}-\widehat p^{m\leftarrow\T}_{r,c}\right|.
\end{align*}
여기서 $c$는 해당 승식의 조합, $C_{r,m}$은 조합 수, $\epsilon$은 자료를 보기
전에 정한 수치 안정화 상수다. $\operatorname{JS}$에는 자연로그를 사용하고
$0\log 0=0$으로 둔다. $\operatorname{MAE}$는 경주 안의 조합별 확률오차를
동일 가중한 값이다. $R^2$는 직관적인 보조 통계로 제시하되, 극단적으로 작은
확률이 많고 한 경주가 조합 수만큼 반복된다는 점 때문에 유일한 적합도 지표로
쓰지 않는다. 두 분포의 상대질량 차이를 보기 위한 centered log-ratio 거리,
상관계수와 cosine similarity도 기술통계로 제시한다.

보정 여부는
\begin{equation}
  \log p^m_{r,c}
  =\alpha_m+\beta_m
   \log\widehat p^{m\leftarrow\T}_{r,c}
   +\delta_{m,n_r}+u_{r,c,m}
  \label{eq:calibration}
\end{equation}
으로 점검한다. 표준오차는 최소한 경주 단위로 군집화한다. 조합 수가 많은 경주가
추정치를 지배하지 않도록 조합가중 결과와 경주균등가중 결과를 나란히 제시한다.

\subsection{기준모형과 위약검정}

삼쌍승식의 높은 차원이 자동으로 높은 설명력을 만드는지 확인하기 위해 다음
기준을 미리 고정한다.

\begin{enumerate}
  \item 단승 정규화 가격으로 식 \eqref{eq:harville}을 만든다.
  \item 같은 경주 안에서 삼쌍승 상태가격을 무작위 순열하여 말의 정체성을 끊는다.
  \item 출전두수가 같은 다른 경주의 삼쌍승 분포를 배정하여 경주 공통정보를 끊는다.
  \item 모든 가능한 삼쌍승 상태에 균등 질량을 주어 인기 구조 자체의 설명력을 제거한다.
\end{enumerate}

모든 비교는 같은 표본, 같은 정규화, 같은 손실함수에서 수행한다. 주 분석의
우위 판정은 하나의 전체표본 $R^2$가 아니라 경주별 손실 차이의 분포와 연도·
경마장 밖 검증에서의 재현 여부에 근거한다.

\subsection{순서정보 가설의 검정}

P3은 같은 경주에서 쌍승식과 복승식의 Harville 대비 개선폭을 짝지어 검정한다.
손실함수 $L$에 대해 삼쌍승식 주변가격의 개선폭을
\[
  \Delta^L_{r,m}
  =L\!\left(p^m_r,A^m_rq^H_r\right)
   -L\!\left(p^m_r,\widehat p^{m\leftarrow\T}_r\right)
\]
로 정의한다. 양의 값은 삼쌍승식 주변가격이 Harville 기준보다 실제 가격에 더
가깝다는 뜻이다. 순서를 구분하는 쌍승식과 순서를 제거하는 복승식의 차이는
\[
  D^L_r=\Delta^L_{r,\Eact}-\Delta^L_{r,\Q}
\]
로 측정한다. 1차 검정은 범위가 승식 차원에 의존하지 않는 TV 거리와
Jensen--Shannon 발산을 사용하고, 동일 경주를 한 쌍으로 한 평균·중앙값과
경주 단위 부트스트랩 신뢰구간을 보고한다. $D^L_r$의 부호와 크기는 순서정보
가설을 검정하지만, 별도 풀의 유동성과 참가자 구성이 함께 다를 수 있으므로
순서 구분의 인과효과로 해석하지 않는다.
또한 이 비교는 순서정보에 관한 축약형 검정이지, 확률가중 함수와 조건부 확률을
명시한 복합복권 비축약의 구조적 검정이 아니다. 후자는 다음 절의 행동모형에서
별도로 수행한다.

\subsection{표본구성, 이질성과 강건성}

가격 정합성 결과에 앞서 Panel A의 clean 경주와 나머지 상한 경주의 출전두수,
연도, 경마장, 상한 조합 수와 상한 직전 비검열 배당의 꼬리분포를 비교한다.
교차풀 오차는 출전두수, 연도, 경마장, 주말 여부, 조합의 인기구간별로 나눈다.
일부 승식의 판매가 중단된 경주와 적중자가 없어 배당률 표시 규칙이 달라진
조합은 별도 표본으로 분리한다.

경주별 통계는 중앙값과 사분위범위를 기본으로 보고하고, 보조적 $R^2$가 0.8을
넘는 경주의 비율과 식 \eqref{eq:calibration}의 기울기가 $[0.9,1.1]$에 드는
경주의 비율을 함께 제시한다. 표준오차는 경주 군집을 기본으로 하고 날짜 수준
공통충격을 허용한 이중 군집 결과를 강건성 분석으로 보고한다.

마지막으로 실현 착순에 대한 로그점수 비교를 보조 분석으로 수행한다. 이는
삼쌍승 가격과 단승--Harville 가격 중 어느 쪽이 실제 순위를 더 잘 예측하는지
묻는다. 내부 가격 일관성과 외부 예측 정확도는 서로 다른 가설이므로 같은 표의
한 계수로 합치지 않는다.

강한 주결론은 Panel A의 점추정과 Panel B의 전체표본 식별구간이 같은 방향을
지지할 때만 내린다. Panel A의 기준모형 대비 개선이 Panel B의 경계에서는 0을
포함하면 `clean 표본에서만 확인`으로 제한하고, 두 패널의 방향이 충돌하면
상한 또는 표본선택 때문에 전체표본 결론이 식별되지 않는다고 판정한다.
